\section{Theorie}

Die Magnetisierung $\vec{M}$ eines Festkörpers ist ein Maß dafür, wie stark ein Material unter Einfluss
eines äußeren Magnetfeldes \(\vec{H}\) magnetisiert wird, also selbst ein Magnetfeld erzeugt.
Sie ist definiert durch:

\begin{equation}
    \vec{M} = \chi_m \vec{H} = (\mu_r - 1) \vec{H}
    \label{eq:Magnetisierung}
\end{equation}

Hierbei ist $\chi_m$ die magnetische Suszeptibilität und $\mu_r$ die relative Permeabilität des Materials.
Die magnetische Suszeptibilität gibt an, wie stark ein Material auf ein äußeres Magnetfeld reagiert. 
Die relative Permeabilität beschreibt die magnetische Durchlässigkeit eines Festkörpers.
Beide Größen sind dimensionslose Materialkonstanten. Für die magnetische Flussdichte in einem Material gild:

\begin{equation}
    \vec{B} = \mu_0 (\vec{H} + \vec{M}) = \mu_0 \mu_r \vec{H}
    \label{eq:Flussdichte}
\end{equation}

Dabei ist $\mu_0$ die magnetische Feldkonstante (Permeabilität des Vakuums) mit dem Wert
$\mu_0 = 1.25663706127 \cdot 10^{-6} \frac{N}{A^2}$ \cite{VacuumPerm}.\\
Über die Permeabilität werden Materalien in drei Klassen eingeteilt:
\begin{itemize}
    \item Diamagnetika mit $\mu_r \leq 1$: \\
    In Stoffen, welche keine internen magnetische Dipole (durch bspw. Spins oder Bahndrehimpulse haben),
    wird durch ein äußeres Magnetfeld ein schwaches Magnetfeld entgegengesetzt zum äußeren Feld induziert(gemäß der Lenz'schen Regel).
    Durch dieses innere induzierte Magnetfeld wird das äußere Feld abgeschwächt.
    \item Paramagnetika mit $\mu_r \geq 1$:\\
    In Stoffen welche keine vollen Elektronenkonfiguration besitzen, gibt es durch ungepaarte Elektronen und Bahndrehimpulse magnetische Dipole.
    Diese Dipole richten sich im äußeren Magnetfeld in die selbe Richtung aus und verstärken dieses somit.
    \item Ferromagnetika mit $\mu_r \gg 1$:\\
    Ferromagnete besitzen starke interne magnetische Dipole, die sich in einem äußeren Magnetfeld gegenseitig ausrichten und somit ein starkes Gesamtmagnetfeld erzeugen.
    Diese Wechselwirkung der Dipole untereinander führt (bis zur Curie-Weiss-Temperatur) zu einer nicht rein statistischen Ausrichtung
    von diesen und damit zu einer permanenten Magnetisierung des Materials auch ohne äußeres Magnetfeld. \\
    Dabei können sich die Dipolmomente in sogenannten Weiss'schen Bezirken (Domänen) spontan parallel ausrichten. Da diese Domänen
    jedoch in verschiedene Richtungen zeigen können, ist die Gesamtmagnetisierung des Materials ohne äußeres Feld oft gering bis null.
    Der Übergangsbereich zwischen Weiss'schen Bezirken wird als Bloch-Wand bezeichnet. \\
    Durch anlegen eines äußeren Magnetfeldes richten sich die Domänen zunehmend in Richtung des äußeren Feldes aus und die Bloch-Wände
    verschieben sich. Da nun der Zusammenhang zwischen Magnetisierung und äußeren Feld durch Störstellen im Festkörper nicht mehr linear ist, entsteht eine Hysteresekurve.
    Die Fläche innerhalb der Hysteresekurve entspricht der Energie, welche für die Verschiebung und Ausrichtung der Bloch-Wände und Weiss'schen Bezierke aufgebracht werden muss.
    Ab einer bestimmten Feldstärke, der Sättigungsfeldstärke $H_c$, sind alle Domänen ausgerichtet und die Magnetisierung ändert sich nicht mehr mit zunehmendem äußeren Feld.
    Bei Abnahme des äußeren Feldes wird der Effekt umgekehrt, jedoch bleibt eine Restmagnetisierung, die sogenannte Remanenz $B_r$ bestehen.
    Wird ein hinreichend starkes Feld in entgegengesetzter Richtung angelegt, so wird das Material entmagnetiseiert. Die hierfür notwendige Feldstärke wird als Koerzitivfeldstärke $H_c$ bezeichnet. 

\end{itemize}
Die Umorientierung der Domänen ist mit Energieverlust verbunden, welcher als Wärme abgegeben wird.
Die pro Messeinheit m dissipierte Energie während eines vollständigen Zyklus der Hysteresekurve lässt sich durch Integration
über die Fläche der Hysteresekurve bestimmen:
\begin{equation}
    \xi_{verlust} = \frac{1}{\rho}\int B dH
    \label{eq:Energieverlust}
\end{equation}